\begin{exo}
	Dans un magasin de meubles, le stock est composé de deux types de meubles :
	\begin{itemize}
		\item $60$\% de canapés et le reste est constitué de fauteuils.
		\item Parmi les fauteuils, $20$\% sont en cuir.
		\item La proportion totale de meubles en cuir (canapés ou fauteuils) dans le magasin est de $44$\%.
	\end{itemize}
	Un client se présente et choisit un meuble au hasard. On considère les événements suivants :
	\begin{itemize}
		\item $F$ : \og{} le meuble choisi est un fauteuil \fg{} ;
		\item $C$ : \og{} le meuble choisi est en cuir \fg{}.
	\end{itemize}

	\begin{enumerate}
		\item Donner, à l'aide des données de l'énoncé, les valeurs de $P(F)$, $P(C)$ et $P_{F}(C)$.
		\item Déterminer la probabilité $P(F \cap C)$ que le meuble choisi soit un fauteuil et qu'il soit en cuir.
		\item Rappeler la formule reliant $P_{C}(F)$, $P(F \cap C)$ et $P(C)$.
		\item À l'aide de cette formule, calculer la probabilité que le meuble choisi soit un fauteuil sachant qu'il est en cuir.
	\end{enumerate}
\end{exo}
